\section{Automatización}
Para automatizar la generación de la geometría se desarrolló una extension de Blender. Para ello, los
grupos de nodos de geometría y los materiales se marcaron como assets y se compilaron en un solo archivo
\textit{.blend} que sirve como fuente de referencia para la extension.
\medskip

La extensión registra tres propiedades para la escena:
\begin{itemize}
    \item \textit{Yolo Settings}: almacena la ruta de archivo del manifesto generado por el componente de
    extracción y los contenidos del mismo.
    \item \textit{Yolo Graph}: contiene un grafo que almacena como nodos la información de cada capa
    de la red, construido a través de los datos importados y deserializados por \textit{Yolo Settings}.
    Su propósito es mantener una referencia entre los datos de una capa y el objeto en la escena
    utilizado para su visualización.
    \item \textit{Yolo Visualization Settings}: almacena variables utilizadas para el control general de
    la visualización, como colores de capas, la escala de los objetos y su espaciado, etc. Su propósito
    es controlar la configuración inicial de la visualización al manipular múltiples objetos de la escena
    al mismo tiempo.
\end{itemize}

La extension permite importar a Blender el manifesto de los datos extraídos (\textit{.json}), almacenando su
información en \textit{Yolo Settings}. Genera también un menú en la vista de la escena para
controlar los parametros de la visualización (\textit{Yolo Visualization Settings}), y un botón
para generar la geometría.
\medskip

Los controles de la visualización son:
\begin{itemize}
    \item \textit{Escala}: multiplicador de escala para el tamaño de la visualización. 
    \item \textit{Escala de Profundidad}: multiplicador de escala específico para la profundidad
    de las capas.
    \item \textit{Espaciado}: distancia entre capas.
    \item \textit{Profundidad Máxima}: profunidad máxima de las capas.
    \item \textit{Color y Simple}: controles por tipo de capa para cambiar el color de capas y operaciones
    del mismo tipo. También permite cambiar la visualización al de un volumen simplificado.
    \item \textit{Generar Geometría}: Genera la geometría de la red.
\end{itemize}

Al generar la geometría, la extensión genera los nodos de \textit{Yolo Graph}, cada uno correspondiendo a
una capa de la red, poblandolos con la información correspondiente de la capa en concordancia con el
manifesto almacenado por \textit{Yolo Settings}. A partir de estos datos genera los grupos de nodos
correspondientes con la información de cada capa, y posteriormente instancia la geometría y grupos de nodos 
para la visualización de las conexiones entre capas. Adicionalmente organiza los objetos en colecciones por
tipo de capa u operación.
